% -----------------------------------------------------------------------------
% Results fragment, theme "method": structure and diagnostics of the pillar indices, the
% rankability rule, the Policy index and the trade inputs. Feeds Section 3 (coverage after
% harmonisation, trade inputs), Section 4 (composition, weights, imputation, diagnostics,
% rankability), Section 8 (Policy index coverage) and Appendix A. Every number is in
% paper/results/method.json and is produced, with the five figures and the four tables
% (tab:method_rankability, tab:method_imputation, tab:method_policy; tables/method_hs6_baskets.tex
% is still written but not input here, since tab:hs6_counts of the misc theme, input in Section 3,
% carries the same counts with trade values), by
% scripts/paper/fig_method.R from the final rerun of 12 October 2026
% (data/processed/outputs/index_outputs.rds, the Comtrade dyad and energy-trade files and the
% HS6 master). "Pair" is a technology-stage pair (t,s) and "cell" a country-technology-stage
% cell (i,t,s); the sample is the eleven headline technologies unless stated otherwise.
% -----------------------------------------------------------------------------

\paragraph{Coverage after harmonisation.} Once names have been mapped to one canonical name
per code (Section~\ref{subsec:data_countries}), the headline pillars cover 219 economies, each
with a distinct ISO3 code. The universe is the same for the two pillars: every economy with an
$\ES$ score in a headline technology also has an $\EO$ score, and conversely. Classified by the
Correlates of War state-system list, 195 of the 219 are sovereign states (the 193 members of the
United Nations, Kosovo and Taiwan) and 24 are dependent or overseas territories, the special
administrative regions of Hong Kong and Macau, and the West Bank and Gaza. Taiwan, which the
World Bank's country table does not list, is carried as an economy of its own throughout,
including the trade themes and the bilateral trade matrix; it has an $\ES$ value in 28 headline
cells and an $\EO$ value in 33. The two pillars score 6,232 and 7,131 country--technology--stage
cells in 30 and 33 headline pairs, and every $\ES$ cell has an $\EO$ value. The grids of both
pillars span almost the whole universe at every stage: 219, 215 and 219 economies at the
upstream, midstream and downstream stages in $\ES$, and 219, 216 and 219 in $\EO$. That width is
set by the most widely distributed theme of each pair rather than by the sources that measure an
economy's own position, a point taken up below. Territories are kept as separate units, but
their scores rest more heavily on imputation: imputed components carry on average 0.48 of the
weight of a territory's headline $\ES$ cell, against 0.29 for a sovereign state, and 0.67 against
0.40 in $\EO$ (the imputed weight share $\iota_{its}$ is defined in
Section~\ref{subsec:norm_diagnostics}). The critical-minerals trade variables use the 2024
Comtrade extracts, with 130 reporters for imports, 128 for exports and 130 for total exports;
the 2025 extracts staged with them have three reporters each, below the floor of 50 in
rule~\eqref{eq:data_resolve}, and were skipped.

\paragraph{Which cells are ranked.} Every cell the sources reach is scored, and every score
enters the decompositions, but a cell enters a cross-country ranking, a top-$N$ list or the
Strategic Index only if it is \emph{rankable}: the economy is a sovereign state and at most
$\tau=0.5$ of the cell's index weight rests on imputed components, $\iota_{its}\le\tau$
(\texttt{config/rankability.yml}). Imputed components here include the zero fills that record
absence from a register, not only the cross-country fills. Figure~\ref{fig:method_rankability} and
Table~\ref{tab:method_rankability} show what the rule keeps. Of the 6,232 scored headline $\ES$
cells, 4,163 (67 percent) are rankable; 607 belong to territories and 1,462 to sovereign states
whose imputed share exceeds the threshold. In $\EO$, 4,301 of 7,131 cells (60 percent) are
rankable, with 718 territory cells and 2,112 sovereign cells over the threshold. Counted over
sovereign cells only, the rule keeps 74 percent in $\ES$ and 67 percent in $\EO$. Its bite
follows the imputation burden described below. In $\ES$ it is heaviest upstream, where only 52
percent of the cells are rankable against 70 percent at the midstream stage and 80 percent
downstream, because the reserve and production registers that dominate upstream $\ES$ cover a
minority of economies and record the others through the zero rule. In $\EO$ it is heaviest at the
midstream and downstream stages (56 and 55 percent, against 70 percent upstream), where the
technology-demand terms are imputed for most economies. The narrowest pairs still keep a usable
comparison set: Green Hydrogen midstream keeps 40 rankable $\EO$ cells (19 percent of those
scored), and Nuclear downstream, the narrowest $\ES$ pair, keeps 68 (32 percent). Every one of
the 195 sovereign states keeps at least one rankable $\ES$ cell, but 17 keep no rankable $\EO$
cell: Liechtenstein and Monaco, and fifteen small island and low-income states (Djibouti,
Dominica, Eritrea, Haiti, Kiribati, the Marshall Islands, Micronesia, Nauru, North Korea, Palau,
S\~ao Tom\'e and Pr\'incipe, Somalia, St Kitts and Nevis, St Vincent and the Grenadines and
Tuvalu). Their only rankable $\EO$ cells were upstream, in pairs without a traded-goods basket; the
upstream baskets added in the final pass bring a Trade category to four of those pairs, whose
export-size and comparative-advantage terms are imputed for economies absent from the Comtrade
reporter file, and this lifts their imputed weight share just over the threshold. Taiwan keeps 19 rankable $\ES$ cells and 18 rankable $\EO$ cells.

The rule is not neutral with respect to the level of the indices. Rankable sovereign cells score
higher on average than sovereign cells over the threshold in every pillar and stage, most
markedly in midstream $\ES$ (0.46 against 0.30) and midstream $\EO$ (0.49 against 0.37). The
reason is mechanical. The cells over the threshold are mostly those of economies absent from the
registers, and the zero rule scores that absence at the bottom of the comparison set, so the
excluded cells lean insecure and less opportune. In the Strategic Index the rule therefore
removes cells with a high measured need: of the 5,854 candidate cells of the policy-scored
economies that have both pillar values in a reported sector, 3,154 (54 percent) are kept, in 174
of 203 economies. The 2,700 cells removed have a mean security need $1-\ES$ of 0.68 and a mean
$\EO$ of 0.44, against 0.58 and 0.49 for the cells kept. The 29 economies that lose every cell
are nine territories and special administrative regions (among them Hong Kong, Macau and the
West Bank and Gaza) and 20 sovereign states: the 17 without a rankable $\EO$ cell, Andorra, the
Central African Republic and Grenada. The rule should be read as a
statement about what the data can support rather than as a correction of the index: an economy's
excluded cells keep their values, and the ranking of the cells that remain is unaffected by their
removal, since the pillars are built before the rule is applied.

\begin{figure}[tbp]
  \centering
  \includegraphics[width=\textwidth]{method_rankability}
  \caption{Rankability of the scored cells by technology--stage pair. Each bar divides the scored
  country--technology--stage cells of a pair into three exclusive groups: rankable cells (the
  economy is a sovereign state and the imputed weight share $\iota_{its}$ is at most
  $\tau=0.5$); cells of sovereign states with $\iota_{its}>0.5$; and cells of territories,
  special administrative regions and the West Bank and Gaza, which are scored but never ranked
  whatever their imputed share. The ``All technologies'' row pools the pairs of the stage, and its
  label gives the rankable share. Sovereignty follows the Correlates of War state-system list (195
  states: the 193 members of the United Nations, Kosovo and Taiwan). ``No ES score'' marks the
  three pairs without an $\ES$ value. Sample: the eleven headline technologies, 6,232 $\ES$ and
  7,131 $\EO$ cells. Source: \texttt{energy\_security\_index} and
  \texttt{economic\_opportunity\_index} in \texttt{data/processed/outputs/index\_outputs.rds};
  \texttt{config/rankability.yml}; \texttt{scripts/paper/fig\_method.R}.}
  \label{fig:method_rankability}
\end{figure}

% Generated by scripts/paper/fig_method.R; do not edit by hand.
\begin{table}[tbp]
  \centering
  \small
  \begin{threeparttable}
    \caption{Rankable cells by pillar and stage, headline technologies}\label{tab:method_rankability}
    \begin{tabular}{@{}l *{4}{r} S[table-format=1.2] r@{}}
      \toprule
      & {Scored} & {Rankable} & \multicolumn{2}{c}{Not rankable} & {Share} & {Economies} \\
      \cmidrule(lr){4-5}
      & {cells} & {cells} & {Territory} & {$\iota_{its}>0.5$} & {rankable} & {ranked} \\
      \midrule
      \multicolumn{7}{@{}l}{\emph{Energy Security}} \\
      \quad Upstream & 2,181 & 1,124 & 231 & 826 & 0.52 & 195 \\
      \quad Midstream & 2,002 & 1,397 & 169 & 436 & 0.70 & 192 \\
      \quad Downstream & 2,049 & 1,642 & 207 & 200 & 0.80 & 193 \\
      \quad All stages & 6,232 & 4,163 & 607 & 1,462 & 0.67 & 195 \\
      \multicolumn{7}{@{}l}{\emph{Economic Opportunity}} \\
      \quad Upstream & 2,394 & 1,681 & 250 & 463 & 0.70 & 178 \\
      \quad Midstream & 2,346 & 1,307 & 221 & 818 & 0.56 & 166 \\
      \quad Downstream & 2,391 & 1,313 & 247 & 831 & 0.55 & 169 \\
      \quad All stages & 7,131 & 4,301 & 718 & 2,112 & 0.60 & 178 \\
      \bottomrule
    \end{tabular}
    \begin{tablenotes}[flushleft]\footnotesize
      \item \emph{Notes:} Country--technology--stage cells of the eleven headline technologies
      with a pillar value. A cell is rankable when the economy is a sovereign state (a Correlates
      of War state-system code) and at most $\tau=0.5$ of its index weight rests on imputed components, $\iota_{its}\le\tau$
      (\texttt{config/rankability.yml}). Reasons are exclusive: a territory is counted as a territory
      whatever its imputed share; the next column counts sovereign cells over the threshold.
      Economies ranked: economies with at least one rankable cell. Source:
      \texttt{data/processed/outputs/index\_outputs.rds}; \texttt{scripts/paper/fig\_method.R}.
    \end{tablenotes}
  \end{threeparttable}
\end{table}


\paragraph{Composition differs across pairs, not within them.}
Figure~\ref{fig:method_category_sets} shows which categories enter each pillar in each of the 33
headline technology--stage pairs. Two features stand out. First, within every pair all scored
economies share one category set. This holds in all 30 pairs with an $\ES$ score and in all 33
with an $\EO$ score. It is the outcome anticipated in Section~\ref{subsec:norm_pillar}: every
theme except the levelised-cost theme carries a missing-data rule, so a category observed for
one economy of a pair is completed for all of them. Cross-country comparisons within a pair are
therefore comparisons of one construct. Second, composition varies widely across pairs. The 30
$\ES$ pairs use 16 distinct category sets and the 33 $\EO$ pairs use 18. No pair observes all ten
categories of either pillar. The number of observed categories ranges from one to seven in $\ES$
(median three) and from two to seven in $\EO$ (median six). Counted over cells rather than pairs
the medians are the same; 21 percent of $\ES$ cells (1,279 of 6,232) are scored on a single
category, whereas no $\EO$ cell has fewer than two.

The variation is mainly structural: it follows the stage at which each category is defined. The
median $\ES$ pair has 5.5 categories upstream, 1.5 at the midstream stage and two downstream.
Only Trade, Foreign Dependency and Investment exist at the midstream stage, and only Consumption,
Energy Access, Investment and, for the three renewable technologies with a resource-potential
series, Reserves exist downstream. In seven of the 30 $\ES$ pairs the pillar consists of a single
category, so $\ES_{its}$ simply is that category's score: Trade in the coal, gas, nuclear and oil
midstream pairs; Foreign Dependency in Green Hydrogen midstream; Investment in Electric Vehicles
downstream; and Consumption in Green Hydrogen downstream. Nuclear upstream, which the uranium
data of the OECD Nuclear Energy Agency and the International Atomic Energy Agency now inform, is
scored on Production, Reserves, Trade and Energy Prices; its Trade category, like that of the
Electric Vehicles, Solar and Wind upstream pairs in both pillars, comes from the upstream trade
baskets described below. The $\EO$ pillar is broader in every pair. Its
narrowest pair, Electric Grid downstream, has only Cost Competitiveness and Technological
Readiness, and the second of these is a technology-level constant. The figure also marks the
categories whose score is identical for every economy in a pair. These are Energy Prices in the
ten upstream pairs that have a price series (in both pillars), Technological Readiness in all 33
$\EO$ pairs, and Technology Demand in five $\EO$ pairs (Batteries upstream and midstream, and
Geothermal at all three stages), where no regional demand outlook separates economies. Such a
category cannot rank economies within its pair. By equation~\eqref{eq:affine_constant}, it shifts
the pair's level and scales down its dispersion.

\begin{figure}[tbp]
  \centering
  \includegraphics[width=\textwidth]{method_category_sets}
  \caption{Category composition of the pillars by technology--stage pair. Each column is one of
  the 33 headline technology--stage pairs, grouped by stage. Each row is a category of the
  Energy Security pillar (upper block) or the Economic Opportunity pillar (lower block), ordered
  by nominal weight as in Table~\ref{tab:category_weights}. Dark blue: the category is observed
  in the pair and its score varies across economies. Light blue: the category is observed but
  takes the same score for every economy of the pair (a world- or technology-level input).
  Gray: the category is not observed in the pair. The number above each column is the number
  of observed categories, which is the same for every economy of the pair. ``Pair not scored''
  marks the three pairs without an $\ES$ value. Sample: the country--technology--stage cells of
  the eleven headline technologies that have a pillar value (6,232 for $\ES$, 7,131 for $\EO$).
  Source: category contributions in \texttt{data/processed/outputs/index\_outputs.rds};
  \texttt{scripts/paper/fig\_method.R}.}
  \label{fig:method_category_sets}
\end{figure}

\paragraph{Effective weights.} Figure~\ref{fig:method_weights} compares the nominal shares of
Table~\ref{tab:category_weights} with the effective weights $\omega_{kits}$ of
equation~\eqref{eq:effective_weight}, averaged over the scored headline cells of each stage and
set to zero where a category is not observed. The category set is common within a pair, so the
effective weight is a property of the pair: averaging over cells instead of pairs changes the
stage means by at most 0.032. The available-case rule hands the nominal weight of the categories
absent from a pair to those present, in proportion to their nominal weights. Because the median
pair observes only three of the ten categories in $\ES$ and five in $\EO$, the effective weights
depart a long way from the nominal ones. The mean absolute gap between effective and nominal
share is 0.11 in $\ES$ and 0.08 in $\EO$, and the largest gap is 0.49.

In $\ES$, Trade has a nominal share of 0.095 but a mean weight of 0.59 at the midstream stage,
and Consumption (nominal share 0.119) carries 0.49 downstream. Upstream the weight is spread more
evenly, led by Reserves (0.20), Trade (0.18) and Production (0.18); Energy Prices, nominally the
smallest $\ES$ category (0.048), carries 0.10. In $\EO$ the redistribution favours the categories
observed at every stage: Technology Demand (0.20 to 0.25, against a nominal 0.133), Cost
Competitiveness (0.19 to 0.24, against 0.100) and Technological Readiness (0.12 to 0.16, against
0.067). The
nominal weights are therefore a poor guide to the composition of any particular cell. This is the
gap between nominal weights and effective importance that \citet{paruolo2013ratings} and
\citet{becker2017weights} emphasise, arising here from which categories are observed rather than
from correlations among indicators.

Inputs that are identical for every economy of a pair take a large share of the $\EO$ weight.
Categories that are constant within the pair (Energy Prices, Technological Readiness and the five
constant Technology Demand pairs) take a mean of 0.34 of the weight of an upstream $\EO$ cell,
0.22 at the midstream stage and 0.18 downstream. Adding the constant components inside categories
that do vary, namely the world export concentration in Trade and the addressable-market,
global-demand and overcapacity terms of Technology Demand, each of which takes one value for
every economy of the pair whether observed or filled, raises these shares to 0.46, 0.39 and
0.27. The corresponding component-level shares in $\ES$ are 0.16, 0.20 and 0. The midstream
$\ES$ figure is the world export concentration term of Trade, which carries a mean weight of
0.20 in nine of the ten midstream pairs; since the export-concentration index is ranked across
the cells of the eleven headline technologies, it sets a pair-level offset rather than a
country-level signal. These shares are a product of the construction, not of the data. They move
pillar levels across pairs without reordering economies within a pair, so they matter for the
within-country comparisons of the Strategic Index (Section~\ref{sec:strategic}) and not for the
within-pair rankings.

\begin{figure}[tbp]
  \centering
  \includegraphics[width=\textwidth]{method_weights}
  \caption{Nominal and effective category weights. Black ticks: nominal share
  $w_k/\sum_{k'}w_{k'}$ over the ten categories of the pillar (\texttt{config/weights.yml};
  Table~\ref{tab:category_weights}). Dots: mean effective weight $\omega_{kits}$
  (equation~\eqref{eq:effective_weight}) over the scored cells of the eleven headline
  technologies at each stage, with $\omega_{kits}=0$ where the category is not observed. A dot
  at zero marks a category that is never observed at that stage. The effective weights of a
  stage, like the nominal shares, sum to one over the categories. A dot to the right of the
  tick means that the available-case rule gives the category more weight than its nominal share.
  Cells: 2,181, 2,002 and 2,049 for $\ES$ and 2,394, 2,346 and 2,391 for $\EO$ (upstream,
  midstream, downstream). Source: category contributions in
  \texttt{data/processed/outputs/index\_outputs.rds}; \texttt{scripts/paper/fig\_method.R}.}
  \label{fig:method_weights}
\end{figure}

\paragraph{Imputation burden.} Figure~\ref{fig:method_imputation} and
Table~\ref{tab:method_imputation} report the mean imputed share of each category by stage and the
mean imputed weight share $\iota_{its}$ of equation~\eqref{eq:imputed_share}. On average, imputed
components carry 31 percent of the weight of an $\ES$ cell and 43 percent of the weight of an
$\EO$ cell (medians 0.26 and 0.44), and 28 percent of $\ES$ cells and 39 percent of $\EO$ cells
exceed the rankability threshold of one half. In $\ES$ the burden is heaviest upstream, at 0.41
against 0.23 at the midstream stage and 0.27 downstream. Upstream, Reserves is imputed at 0.84 on
average and alone accounts for 0.17 of the stage mean of 0.41. Its values come from the zero rule,
which records absence from the reserve registers (the U.S. Geological Survey's mineral reserves,
the Energy Institute's fuel reserves and the uranium resources of the Red Book) as no reserves.
Energy Imports (0.55), Foreign Dependency (0.50) and Production (0.43) also cover only part of
the upstream grid, and the new upstream Trade category is imputed at 0.19. In $\EO$ the burden is
flat across stages (0.41, 0.44 and 0.43) and is led by
Technology Demand, whose score is imputed at 0.73 to 0.86 on average. Three of its terms cover
only part of the grid. The addressable-market term is a market-level value observed for one
market per pair, so it counts as imputed for every other economy; it supplies 0.37 to 0.48 of the
category's mean imputed share. The regional demand outlook, the category's only country-varying
term, is observed for a median of 12 economies per pair and supplies 0.22 to 0.33. The
global-demand term is a world-level series recorded for 108 economies and filled with the same
value for the others: it is identical for every economy, yet it counts as imputed in about half
of the cells and supplies 0.11 to 0.13. Technology Demand contributes between 0.16 and 0.19 to
each stage mean of $\iota_{its}$. Investment, which enters both pillars, is heavily imputed at
every stage of each (0.64 to 0.83): where the category is defined, its components are observed
for only a minority of the economies in the pair's grid.

Categories built from world-level series show almost no imputation: less than 0.01 for Energy
Prices and at most 0.04 for Technological Readiness. This illustrates the point of
Section~\ref{subsec:norm_diagnostics} that $\iota_{its}$ records how a value was obtained, not
whether it is specific to the economy, and it qualifies the rankability rule, which can admit a
cell whose observed weight is largely world-level. No $\ES$ cell rests entirely on imputed
components. In $\EO$, 24 cells do, all in three small economies: the Isle of Man (14 cells),
Puerto Rico (8) and Liechtenstein (2). Nineteen of the 24 are at the midstream stage. These cells
are kept and flagged but are never ranked. At the other extreme, 2,055 $\ES$ cells (33 percent)
and 407 $\EO$ cells (6 percent) contain no imputed component.

\begin{figure}[tbp]
  \centering
  \includegraphics[width=\textwidth]{method_imputation}
  \caption{Imputation burden by category and stage. Category rows: the mean, over the headline
  cells in which the category is observed, of the share of the category score carried by
  imputed components, $\sum_v\lambda_{vkits}\1[\tilde x_{itsv}\text{ imputed}]$. Pillar row
  (with values): the mean imputed weight share $\iota_{its}$ of
  equation~\eqref{eq:imputed_share} over all scored cells of the stage, which weights the
  category rows by the effective weights. Dashes mark categories that are not observed at a
  stage. A component counts as imputed when a missing-data rule of
  Table~\ref{tab:missing_rules} supplied it, zero fills included; world-level series copied to
  every economy count as observed. All values are given in Table~\ref{tab:method_imputation}.
  Source: variable contributions and pillar diagnostics in
  \texttt{data/processed/outputs/index\_outputs.rds}; \texttt{scripts/paper/fig\_method.R}.}
  \label{fig:method_imputation}
\end{figure}

% Generated by scripts/paper/fig_method.R; do not edit by hand.
\begin{table}[tbp]
  \centering
  \small
  \begin{threeparttable}
    \caption{Mean imputed share by category, stage and pillar, headline technologies}\label{tab:method_imputation}
    \begin{tabular}{@{}l *{6}{S[table-format=1.2]}@{}}
      \toprule
      & \multicolumn{3}{c}{Energy Security} & \multicolumn{3}{c@{}}{Economic Opportunity}\\
      \cmidrule(lr){2-4}\cmidrule(l){5-7}
      {Category} & {Up} & {Mid} & {Down} & {Up} & {Mid} & {Down} \\
      \midrule
      Consumption & {--} & {--} & 0.03 & {--} & {--} & 0.05 \\
      Cost Competitiveness & {n.a.} & {n.a.} & {n.a.} & 0.34 & 0.42 & 0.36 \\
      Energy Access & {--} & {--} & 0.65 & {--} & {--} & 0.66 \\
      Energy Imports & 0.55 & {--} & {--} & {n.a.} & {n.a.} & {n.a.} \\
      Energy Prices & 0.01 & {--} & {--} & 0.01 & {--} & {--} \\
      Foreign Dependency & 0.50 & 0.54 & {--} & {--} & 0.79 & {--} \\
      Investment & 0.83 & 0.78 & 0.64 & 0.83 & 0.78 & 0.66 \\
      Minerals Trade & 0.02 & {--} & {--} & {n.a.} & {n.a.} & {n.a.} \\
      Production & 0.43 & {--} & {--} & 0.68 & {--} & 0.55 \\
      Reserves & 0.84 & {--} & 0.10 & {n.a.} & {n.a.} & {n.a.} \\
      Technological Readiness & {n.a.} & {n.a.} & {n.a.} & 0.04 & 0.02 & 0.03 \\
      Technology Demand & {n.a.} & {n.a.} & {n.a.} & 0.86 & 0.73 & 0.83 \\
      Trade & 0.19 & 0.01 & {--} & 0.26 & 0.15 & {--} \\
      \midrule
      Pillar, $\bar\iota$ & 0.41 & 0.23 & 0.27 & 0.41 & 0.44 & 0.43 \\
      Cells & {2,181} & {2,002} & {2,049} & {2,394} & {2,346} & {2,391} \\
      \bottomrule
    \end{tabular}
    \begin{tablenotes}[flushleft]\footnotesize
      \item \emph{Notes:} Category rows give the mean, over the country-technology-stage cells
      of the eleven headline technologies in which the category is observed, of the share of the
      category score carried by imputed components,
      $\sum_v\lambda_{vkits}\1[\tilde x_{itsv}\text{ imputed}]$. The pillar row gives the mean
      over all scored cells of the stage of the imputed weight share $\iota_{its}$ of
      equation~\eqref{eq:imputed_share}. --: category not observed at that stage; n.a.: category
      not part of the pillar. Source: \texttt{data/processed/outputs/index\_outputs.rds}
      (variable contributions and the pillar diagnostics); \texttt{scripts/paper/fig\_method.R}.
    \end{tablenotes}
  \end{threeparttable}
\end{table}


\paragraph{Coverage, country-specific information, floors and pair status.}
Figure~\ref{fig:method_coverage} plots two counts for each pair. The first is the number of
economies scored. The second is the number of those economies with at least one positively
weighted, non-imputed component that is country-specific, meaning that it takes different
observed values across the economies of the pair. The gap between the two is the number of
economies whose score in the pair rests only on imputed values and on world-level series. In
$\ES$ there is no gap at the midstream and downstream stages: in all 20 of those pairs every
scored economy has country-specific data. Upstream, where the commodity-price series copied to
every economy sets the grid, every pair except Nuclear has a gap, from four economies in Electric
Vehicles to 38 in Coal, where 180 of 218 economies have country-specific $\ES$ data. Nuclear
upstream has none: every one of its 218 economies has a country-specific value once absence from
the uranium registers is recorded as zero production and zero resources. In $\EO$, every pair
except Nuclear upstream has at least one scored economy without country-specific data. The share
with such data is lowest in Electric Grid downstream (143 of 213 economies, 0.67), followed by
the upstream and midstream geothermal pairs (144 of 213 each), Green Hydrogen midstream (148 of
213) and Green Hydrogen upstream (150 of 218). Every economy in
the upstream grids has a country-specific upstream component in at least one pair of each pillar.

The economies without country-specific data share one value within their pair, which shows up
as a mass point: in five $\EO$ pairs (Electric Grid downstream, Green Hydrogen upstream and
midstream, and Geothermal upstream and midstream) between 65 and 70 economies, close to a third of
those scored, sit at a single interior value. None of these tied cells is rankable, so the rule
of the previous paragraph removes them from every ranking. A second kind of mass point is
produced by observed zeros, which the rule does not remove. In $\ES$, 98 of the 109 economies
scored in Green Hydrogen midstream share the pair's minimum value (0.256), and 95 of these cells
are rankable. In Nuclear upstream, 127 of the 218 economies share one value (0.172), 105 of them
rankable; the tie is no longer the pair's minimum, because the uranium Trade category places some
economies below it. Such a tie records a real absence of production or resources, but it gives
many economies the same low security score; with the security-need reading of the Strategic
Index it becomes a uniform, sector-wide need (Section~\ref{sec:strategic}).

All 30 headline pairs with an $\ES$ score and all 33 with an $\EO$ score have status ok; none is
constant and none is unscored. Over every technology, including those outside the headline set,
all 45 $\ES$ pairs and all 40 $\EO$ pairs are ok. Nuclear upstream, which was constant before the
uranium data were added, now varies across economies. Three pairs, Electric Grid downstream and
Geothermal upstream and midstream, have no observed $\ES$ category and so no $\ES$ score. The
Strategic Index is therefore computed over 30 sectors: the degenerate-sector rule of
Section~\ref{subsec:si_sample} removes none, Geothermal midstream is left out by rule, and
Electric Grid downstream and Geothermal upstream have no $\ES$ value.

\begin{figure}[tbp]
  \centering
  \includegraphics[width=\textwidth]{method_coverage}
  \caption{Economies scored and economies with country-specific data, by technology--stage pair.
  Filled dots: the number of economies with a pillar value in the pair. Hollow circles: the
  number of those economies for which at least one positively weighted, non-imputed component is
  country-specific, that is, takes at least two distinct observed values across the economies of
  the pair. World-level series copied to every economy, such as commodity-price volatility, world
  export concentration and technology readiness, do not count as country-specific. The segment
  between the two marks is the number of economies whose value rests only on imputed components
  and world-level series. Where the two counts coincide the marks appear as a ringed dot. ``No ES
  score'' marks the pairs without an $\ES$ value. Every pair has status ok (no pair is constant or
  unscored). The axis starts at 100 economies. Sample: the eleven headline technologies. Source:
  variable contributions and pillar diagnostics in
  \texttt{data/processed/outputs/index\_outputs.rds}; \texttt{scripts/paper/fig\_method.R}.}
  \label{fig:method_coverage}
\end{figure}

\paragraph{The Policy index.} Table~\ref{tab:method_policy} summarises the composite used in
the final run. It has two components, the IEA Policies and Measures database
\citep{iea2026pams} with weight 3 and the Climate Action Tracker ratings \citep{cat2026ratings}
with weight 2. The New Industrial Policy Observatory component \citep{gta2026nipo} is excluded
and survives only in a robustness variant, built from the October 2026 release of the Observatory
(interventions in force on 5 October 2026) in the same way as the pipeline builds the component;
it covers 77 jurisdictions, with the EU members scored individually. The dual-use score is
excluded because it takes one value for every economy of a pair and so cannot rank economies. IEA PAMS is required and CAT is added where it
rates the economy, with the weights renormalised over the components present. IEA PAMS reaches
215 economies, and all 215 receive a Policy index: 65 on both components (39 with a rating of
their own and 26 EU members that take the EU's rating) and 150 on IEA PAMS alone. Every one of
the 57 coalition members is scored, 11 of them on IEA PAMS alone (Israel and Ukraine in the core;
the Democratic Republic of the Congo, the Dominican Republic, Gabon, Ghana, Guatemala, Malaysia,
Namibia, Oman and Qatar among the partners). Requiring both components, the rule of earlier
drafts, leaves 65 economies and is kept as a robustness variant.

The basis on which an economy is scored matters for the within-country comparisons that the
Strategic Index makes. For the 65 economies scored on both components, the within-country
Spearman correlation between the IEA PAMS component alone and the composite, computed over each
economy's headline cells, has a median of 0.48 and an interquartile range of 0.41 to 0.57; it is
below one half for 37 of the 65. Over the cells that enter the Strategic Index the median is
0.44. The two components carry largely independent information: across the cells of these 65
economies their Spearman correlation is $-0.05$. Adding CAT therefore reorders an economy's
sectors substantially, and the 109 economies of the Strategic Index scored on IEA PAMS alone,
which hold 1,488 of its 3,154 cells, are ordered on a different basis from the 65 that hold the
other 1,666. The comparison is within economies, so the two groups are not mixed in any single
ranking of the policy term, but the difference in basis should be kept in mind when the
Strategic Index of a member scored on IEA PAMS alone is compared with that of a member scored on
both.

% Generated by scripts/paper/fig_method.R; do not edit by hand.
\begin{table}[tbp]
  \centering
  \small
  \begin{threeparttable}
    \caption{Components of the Policy index}\label{tab:method_policy}
    \begin{tabular}{@{}llrl@{}}
      \toprule
      Component & Weight & Economies & In the composite \\
      \midrule
      IEA Policies and Measures (PAMS) & 3 & 215 & yes, required \\
      Climate Action Tracker rating & 2 & 65 & where rated (39 own, 26 via the EU) \\
      GTA New Industrial Policy Observatory & -- & 77 & no (robustness variant, October 2026) \\
      Dual-use policy score & -- & 217 & no (constant within pairs) \\
      \midrule
      Economies with a Policy index & & 215 & \\
      \quad on IEA PAMS and CAT & & 65 & \\
      \quad on IEA PAMS only & & 150 & CAT weight renormalised away \\
      Coalition members with a Policy index & & 57 & of 57 (11 on IEA PAMS only) \\
      \bottomrule
    \end{tabular}
    \begin{tablenotes}[flushleft]\footnotesize
      \item \emph{Notes:} Economies: economies with at least one value of the component in the
      eleven headline technologies. The Policy index of a cell is the weighted mean of the
      composite components present in it, with the weights renormalised over those components.
      IEA PAMS is required: an economy is scored when it has an IEA PAMS value in at least one
      headline cell, and the Climate Action Tracker rating is added where the economy has one.
      EU member states without a rating of their own take the EU's rating. Requiring both
      components (\texttt{OPSI\_POLICY\_REQUIRE\_ALL=1}, a robustness variant) leaves the 65 economies scored on both. For those 65, the within-country Spearman correlation between the IEA PAMS component alone and the composite over their headline cells has median 0.48 (interquartile range 0.41--0.57), so the basis on which an economy is scored affects the ordering of its sectors.
      The NIPO component enters only a robustness variant, built from the October 2026 NIPO release (interventions in force on 5 October 2026; 77 jurisdictions, EU members individually); the dual-use score takes one value for every economy of a technology--stage pair, so it cannot rank economies.
      Source: \texttt{policy\_component\_tbl}, \texttt{policy\_eligibility}, \texttt{policy\_index} and
      \texttt{coalition\_members} in \texttt{data/processed/outputs/index\_outputs.rds};
      \texttt{scripts/paper/fig\_method.R}.
    \end{tablenotes}
  \end{threeparttable}
\end{table}


\paragraph{Trade inputs: baskets, bilateral flows and conduits.}
Table~\ref{tab:hs6_counts} in Section~\ref{sec:data} gives the HS6 baskets after the basket review
of October 2026.
The baskets hold 259 distinct codes, 242 of the current HS revision and 17 legacy codes that
bridge the HS 1992 nomenclature of the Atlas panels; 105 codes belong to more than one basket. Of
the 33 headline pairs, 18 upstream and midstream pairs have a basket; only the Geothermal and Green
Hydrogen upstream and midstream pairs have none, and no basket exists for the downstream stage.
Four of the eighteen are upstream baskets added in the final pass, which feed the Trade categories
of both pillars: Electric Vehicles (eight codes: the ores of the battery minerals and the rare
earths of permanent magnets), Wind (three codes: rare earths), Nuclear (two codes: uranium ores
and natural uranium compounds, the heading of most traded ore concentrate, which moved from the
midstream basket) and Solar (one code: silicon
metal). A Semiconductors upstream basket (gallium, germanium and indium) is carried for
description only. The review took 92 rows out of the master,
each with its reason recorded: 47 codes leave every basket and 22 stay in another technology's
basket. It removed general-purpose goods whose main use lies outside the technology (aircraft
engines and building boilers from coal and gas midstream, commodity resins and packaging film
from batteries midstream, general machinery from oil and gas midstream, telecommunication cables
from the grid, solar and wind baskets) and two items with direct consequences for the results.
Electricity (HS 271600) is no longer in the Electric Grid upstream basket: the commodity the grid
carries is not an input to building it, and with it in the basket, net power exporters ranked as
grid suppliers. The pair's basket is now the copper and aluminium material inputs, seven codes
from copper ores, blister and cathodes to bauxite, alumina and unwrought aluminium. Switching and
routing apparatus (HS 851762), which is telecommunication and data-networking equipment, is no
longer in Electric Grid midstream.

Bilateral flows come from an all-reporter pull of Comtrade dyads \citep{uncomtrade} for 2020 and
2024, in which 173 economies report in at least one of the two years (168 in 2020; 142 for
exports and 143 for imports in 2024) and 208 appear as reporter or partner; Taiwan, which Comtrade
records as ``Other Asia, n.e.s.'', is mapped to its own code. Each exporter--importer flow is the
importer's declaration where it exists and otherwise the exporter's declaration scaled by a
CIF/FOB factor of 1.06. Of the 701,214 exporter--importer--technology--stage cells observed in
2024, 22 percent come from mirror declarations; they are 13 percent of the cells with positive
trade and 4 percent of the value, because large traders report. Thirty importers that reported in
2020 did not report in 2024, among them Russia, Iran, the United Arab Emirates and Viet Nam, and
65 of the 208 importers enter the 2024 matrix only through their partners' declarations. Their
imports are therefore measured from the exporters' side, which omits flows from exporters that do
not report either. The country-level trade themes of the pillars read a file of exports and
imports by reporter and code that is rebuilt from the World rows of the same pull. It covers 142
exporting and 143 importing reporters in 2024 and every basket code of the current HS revision
(the thirteen codes absent from it are legacy codes); among others it restores Qatar's exports of
liquefied natural gas, USD 56.9 billion in 2024. Hong Kong is left out of the Comtrade export
measures (export size, revealed comparative advantage and export concentration) because its
exports are mostly re-exports of goods made elsewhere, which the producers' own exports already
count; its Atlas-based variables are kept. In the outbound direct-investment component of the
bilateral indices, positions with 16 conduit financial centres are treated as unobserved: the ten
largest conduits identified by Damgaard, Elkjaer and Johannesen (the Netherlands, Luxembourg,
Hong Kong, the British Virgin Islands, Bermuda, Singapore, the Cayman Islands, Switzerland,
Ireland and Mauritius) and six further offshore centres (the Bahamas, Barbados, the Turks and
Caicos Islands, the Isle of Man, Guernsey and Cyprus).

\paragraph{Exactness of the decompositions and mixing of vintages.} The decompositions of
equations~\eqref{eq:cat_decomp} and~\eqref{eq:var_decomp} hold to machine precision. Over all
7,794 $\ES$ cells and 7,661 $\EO$ cells, covering every technology, the largest absolute
difference between the index and the sum of its category contributions is zero, as it must be,
since the builders compute the index as that sum; the check confirms that the stored
contributions are the ones that produced the index. The largest difference between the index
and the sum of its variable contributions is $1.1\times10^{-16}$, a rounding error. In every
cell the realised category weights sum to one, as do the variable weights
$\omega_{kits}\lambda_{vkits}$. Rebuilding $\iota_{its}$ from the variable table reproduces the
builders' diagnostic to within $1.1\times10^{-16}$. The attributions in
Sections~\ref{sec:energy_security} and~\ref{sec:economic_opportunity} are therefore exact rather
than approximate, and they hold for every scored cell, rankable or not.

The latest-year rule of equation~\eqref{eq:latest_year} mixes vintages within cells. In $\ES$,
60 percent of cells combine components from more than one year, at most three, all between 2023
and 2025 apart from the 2035 horizon of the announced manufacturing-capacity shares. In $\EO$ the
share is 92 percent, and a cell can contain up to six distinct years. Part of the $\EO$ mixing
reflects the 2030 and 2035 projection horizons of the outlook series rather than data vintages.
Counting only years up to 2026, 91 percent of $\EO$ cells still combine more than one vintage,
with at most four years and a median spread of three years among the cells that mix. The extreme
case is the input-cost index of Papua New Guinea, whose latest observation is for 1980 and which enters all 33 of that economy's $\EO$ cells through Cost Competitiveness.
The 2035 manufacturing-capacity shares carry a limitation of their own. The IEA reports the
announced capacity for Europe as a region, so every EU member takes the EU's regional share in the
market-share terms of midstream Foreign Dependency (and of the corresponding $\EO$ variable);
these terms therefore do not separate the EU members from one another, and an EU member's share
of planned capacity is overstated where the EU's capacity is located elsewhere in the Union.
